% Source physical page 2, lower portion. All eight entries checked against scan.
\begingroup
\sffamily\fontsize{8}{10.5}\selectfont
\setlength{\parindent}{0pt}\setlength{\parskip}{4pt}
\noindent{\Large 物理学讲义}\par\vspace{3pt}\hrule\vspace{4pt}
\textbf{已出版}\par
\noindent\begin{minipage}[t]{.49\linewidth}
\textbf{第 1 卷:} J. C. Erdmann, 晶体中的热传导: 理论基础与先进的实验方法. II, 283 页. 1969. DM 20,—\par
\textbf{第 2 卷:} K. Hepp, 重整化理论. III, 215 页. 1969. DM 18,—\par
\textbf{第 3 卷:} A. Martin, 散射理论: 幺正性、解析性与交叉性. IV, 125 页. 1969. DM 14,—\par
\textbf{第 4 卷:} G. Ludwig, 物理理论概念的阐释, 以及通过测量基本定律为量子力学的 Hilbert 空间结构建立公理基础. XI, 469 页. 1970. DM 28,—\par
\textbf{第 5 卷:} M. Schaaf, 正向保时量子力学 Poincaré 群的两个不可约幺正表示的乘积的约化. IV, 120 页. 1970. DM 14,—\par
\textbf{第 6 卷:} 数学与物理中的群表示. V. Bargmann 编. V, 340 页. 1970. DM 24,—\par
\textbf{第 7 卷:} R. Balescu, J. L. Lebowitz, I. Prigogine, P. Résibois, Z. W. Salsburg, 统计物理讲义. V, 181 页. 1971. DM 18,—\par
\textbf{第 8 卷:} 第二届流体动力学数值方法国际会议文集. M. Holt 编. IX, 462 页. 1971. DM 28,—\par
\end{minipage}
\endgroup
